\section{Experimental Setup}
\label{sec:experiments}

\subsection{Implementation}
% \todo{Fill in: framework (PyTorch), pyKT version/commit, hardware (GPU),
% and the public repository URL for the reproducible benchmark.}
All models are implemented in PyTorch on top of the pyKT
toolkit~\cite{liu2022pykt}. Code and configurations are released for
reproducibility.

\subsection{Hyperparameters}
Models are trained with the Adam optimiser and binary cross-entropy loss,
with early stopping (patience~10 epochs, maximum~50) on a validation split
and a maximum sequence length of~200. Hyperparameters are selected by
searching the grid in Table~\ref{tab:hparams}; the final column lists the
values used for the proposed model. \emph{Note:} the final values follow
common practice for the attention baseline and may be re-tuned per dataset.

\begin{table}[t]
\caption{Hyperparameter search grid and final values for the proposed
model. Re-tune per dataset if needed.}
\label{tab:hparams}
\centering
\small
\begin{tabular}{@{}lll@{}}
\toprule
\textbf{Hyperparameter} & \textbf{Search grid} & \textbf{Final} \\
\midrule
Embedding size    & $\{64, 128, 256\}$        & $256$ \\
Hidden size       & $\{64, 128, 256\}$        & $256$ \\
Attention heads   & $\{4, 8\}$                & $8$ \\
Attention blocks  & $\{1, 2\}$                & $2$ \\
Dropout           & $\{0.1, 0.2, 0.3\}$       & $0.2$ \\
Learning rate     & $\{10^{-3}, 10^{-4}\}$    & $10^{-4}$ \\
Batch size        & $\{32, 64\}$              & $32$ \\
\bottomrule
\end{tabular}
\end{table}

\subsection{Evaluation Protocol}
We use 5-fold student-level cross-validation. Metrics are AUC, RMSE, and
accuracy for prediction; and the sequencing-utility metric defined in
Section~\ref{sec:sequencing} for downstream evaluation. We report means
and standard deviations, and test significance with paired tests across
folds together with effect sizes.

\subsection{Research Questions Restated}
\begin{itemize}
  \item \textbf{RQ(a):} Do KC/code-structure-aware models predict mastery
        more accurately than response-only models?
  \item \textbf{RQ(b):} Do accuracy gains translate into better
        next-exercise sequencing?
  \item \textbf{RQ(gen):} How well do findings generalise across datasets
        and relative to the mathematics domain?
\end{itemize}

\subsection{Experimental Hypotheses}
\label{sec:hypotheses}
We state falsifiable hypotheses for each research question, together with
the test used and the decision criterion. Unless otherwise noted, tests
are two-sided at significance level $\alpha = 0.05$, computed over the
five cross-validation folds, and we report an effect size alongside each
$p$-value. Here $\mathrm{AUC}_{\text{full}}$ denotes the AUC of the full
proposed model (response~+~KC~+~code), and $\mathrm{AUC}_{\text{resp}}$
that of the strongest response-only baseline.

\begin{itemize}
  \item \textbf{H1 (prediction; addresses RQ(a)).}
        The full model predicts more accurately than the best
        response-only baseline:
        $\mathrm{AUC}_{\text{full}} > \mathrm{AUC}_{\text{resp}}$.
        \textit{Null} $H1_0$: $\mathrm{AUC}_{\text{full}} \leq
        \mathrm{AUC}_{\text{resp}}$.
        \textit{Test}: paired $t$-test (or Wilcoxon signed-rank) across
        folds on per-fold AUC; effect size Cohen's $d$.

  \item \textbf{H1a (ablation; addresses RQ(a)).}
        Each signal contributes positively: both the $+$KC and the
        $+$code variants improve over the response-only baseline, and the
        combined model is no worse than either single-signal variant.
        \textit{Test}: paired comparisons across folds with
        Holm--Bonferroni correction for multiple comparisons.

  \item \textbf{H2 (sequencing; addresses RQ(b)).}
        The full model yields higher offline sequencing utility
        $U$ than the best response-only baseline:
        $U_{\text{full}} > U_{\text{resp}}$.
        \textit{Null} $H2_0$: $U_{\text{full}} \leq U_{\text{resp}}$.
        \textit{Test}: paired test across folds on per-fold utility;
        effect size reported.

  \item \textbf{H3 (accuracy--utility coupling; core of RQ(b)).}
        Across all models, predictive accuracy and sequencing utility are
        positively correlated:
        $\rho(\mathrm{AUC}, U) > 0$.
        \textit{Null} $H3_0$: $\rho \leq 0$.
        \textit{Test}: Pearson and Spearman correlation over the set of
        models, with a bootstrap confidence interval on $\rho$.
        A weak or non-significant $\rho$ is itself a substantive
        finding: it would indicate that accuracy gains do \emph{not}
        reliably translate into better sequencing.

  \item \textbf{H4 (cross-dataset generalisation; addresses RQ(gen)).}
        A model trained on one programming dataset and tested on another
        retains predictive skill above chance
        ($\mathrm{AUC} > 0.5$), though we expect a drop relative to
        within-dataset performance.

  \item \textbf{H5 (cross-domain contrast; addresses RQ(gen)).}
        The benefit of programming-specific signals (KC~+~code) is larger
        on programming datasets than any transferable benefit observed
        when contrasting with the mathematics domain, where the
        code-structure signal is unavailable.
\end{itemize}

Table~\ref{tab:hyp-map} maps each hypothesis to the results table that
tests it.

\begin{table}[t]
\caption{Mapping of hypotheses to results.}
\label{tab:hyp-map}
\centering
\small
\begin{tabular}{@{}lll@{}}
\toprule
\textbf{Hypothesis} & \textbf{Research question} & \textbf{Evidence in} \\
\midrule
H1, H1a & RQ(a)   & Table~\ref{tab:prediction} \\
H2, H3  & RQ(b)   & Table~\ref{tab:sequencing} \\
H4, H5  & RQ(gen) & Table~\ref{tab:transfer} \\
\bottomrule
\end{tabular}
\end{table}
